\chapter{Хэлэлцүүлэг}
\label{ch:discussion}

%======================================================================
% IMRaD хэлэлцүүлэг (11/07): үр дүнг өмнөх судалгаатай харьцуулах →
% практик утга → хязгаарлалт. Туршилт тус бүрийг тусад нь.
% ӨӨРИЙН ҮГЭЭР ЗАСАХ: энэ бүлэг таны өөрийн тайлбар байх ёстой.
%======================================================================

\section{Дуу таних}
\label{sec:disc-stt}

Whisper-ийн анхны өгүүлэлд олон хэлний Whisper large-v2 загвар Монгол хэлэнд 100\%-иас их WER авсан \cite{radford2023whisper}. Манай Монгол хэлээр нэмэлт сургасан Whisper small загвар 38\% WER авсан. Mengke нар ч Whisper-ийг Монгол өгөгдлөөр сургахад алдаа их буурсныг харуулсан \cite{mengke2025khalkha}. Энэ нь Монгол хэлэнд нэмэлт сургалт чухал гэдгийг дэмжиж байна. Гэхдээ эдгээр өгүүллийн өгөгдөл манайхаас өөр тул тоог шууд харьцуулахгүй.

Gemini 3.5 Transcribe нэмэлт сургалтгүйгээр Монгол хэлээр сургасан Whisper-тэй ойролцоо үгийн алдаатай гарсан. Харин тэмдэгтийн алдаа их, нэг бичлэгт оныг буруу таньсан. Сургуулийн мэдээлэлд он, тоо чухал тул энэ сул тал хариултын чанарт шууд нөлөөлнө.

Хэмжилтэд Gemini илүү гараагүй ч 2026 оны 11-р сард чатад Gemini 3.5 Transcribe-ийг ашиглахаар шийдсэн. Шийдвэрийн шалтгаан нь яриаг текст болгох тусгай үйлчилгээ, компьютерт загвар ажиллуулах шаардлагагүй, үнэгүй квоттой, сургуулийн нэр томьёог жагсаалтаар зааж өгөх боломжтой. Сул тал нь интернэт, API түлхүүр хэрэгтэй, аудио Google руу очдог. Whisper-ийн дүн нь харьцуулалтын суурь болж үлдэнэ.
% ӨӨРИЙН ҮГЭЭР ЗАСАХ: шийдвэрийн шалтгааныг багштай ярилцсан үгээрээ бич
% (багшийн үгийг зохиохгүй).

\section{Хайлт ба эхтэй хариулт}
\label{sec:disc-rag}

Lewis нарын туршилтад векторын хайлт ихэнх ажилд BM25-аас илүү, нэр их агуулсан асуултад BM25 илүү байсан \cite{lewis2020rag}. Манай 20 асуултад E5 нь BM25-аас бага зэрэг илүү гарсан. Хоёуланг нь нэгтгэхэд бүх асуултын зөв хэсэг эхний 3-т орсон. Асуулт цөөн тул аль арга нь илүү гэж дүгнэхгүй. Гэхдээ хоёр аргыг хослуулах нь манай өгөгдөлд муудуулаагүй.

Хайлт сайн байсан ч зөв хариулт 20-оос 13 байв. Алдааны ихэнх нь хариулт сонгох дүрэмд гарсан. Дүрэм асуултын үгийг эхийн үгтэй шууд тулгадаг. Монгол хэлэнд үгийн төгсгөл олон янзаар хувирдаг тул «хэддүгээр», «индекс код» зэрэг хэллэг эхтэй таарахгүй байна. Lewis нарын системд хэлний загвар олон баримтыг уншиж хариулт бичдэг \cite{lewis2020rag}. Манай систем хариултыг шинээр бичихгүй, эхийн үгээр өгдөг. Үүний ашиг нь хариулт бүр эхтэйгээ шууд холбогдсон байх. Сул тал нь хэллэг бага зэрэг өөр байхад хариулт олдохгүй байх.

Эхгүй 19 асуултын 7-д систем ойролцоо хэсгээр хариулсан. Хэрэглэгчийн хувьд буруу хариулт «мэдэхгүй» гэхээс илүү хор хөнөөлтэй. B өөрчлөлт эхгүй хариултыг цөөлсөн ч зөв хариултыг ч хамт цөөлсөн. Нэг дүрмийг чангатгах нь хоёр талдаа нөлөөлж байна. Энэ асуудлыг шийдэхийн тулд эхгүй асуултыг тусад нь танидаг аргыг судлах хэрэгтэй.

Lewis нар баримтын санг солиход мэдлэг шинэчлэгдэхийг харуулсан \cite{lewis2020rag}. Манай системд ч сургуулийн мэдээлэл өөрчлөгдөхөд загвар сургахгүй, зөвхөн хадгалсан хуудсыг шинэчилнэ. Мэдлэгийн мод нь сангийн аль хэсэг хоосон байгааг харуулдаг. Одоогоор багш, өрөө, хуваарийн мэдээлэл байхгүй тул ресепшний асуултад систем хариулж чадахгүй. Энэ нь хайлтын арга биш, өгөгдлийн дутагдал юм.
% ИШ ТАТАХ: эхгүй асуултын судалгааг (SQuAD 2.0, RAG08) W03-д уншсаны дараа энд холбоно.

\section{Хязгаарлалт}
\label{sec:limits}

\textbf{Туршилт 1 (дуу таних):}
\begin{enumerate}
    \item \textbf{Нэг хүн.} Бүх бичлэг нэг хүнийх. Өөр хүний хоолой, аялга, орчинд үр дүн өөр байж болно.
    \item \textbf{Цөөн бичлэг.} 12 бичлэг, 176 үг. Интервал өргөн бөгөөд бодит тодорхойгүй байдал үүнээс ч их байж болно (\ref{sec:stt-ci}).
    \item \textbf{Хөгжүүлэлтийн багц.} Бичлэгүүдийг өмнө нь системийг хөгжүүлэхэд харсан. Тусдаа шалгалтын багц биш.
    \item \textbf{Нэг удаагийн хэмжилт.} Gemini-д бичлэг бүрийг нэг удаа илгээсэн. Дахин илгээхэд ижил текст гарах эсэхийг шалгаагүй. Нэг бичлэг хязгаарын алдаагаар хасагдсан.
    \item \textbf{Үйлчилгээний хувилбар.} Google загвараа шинэчилж болно. Энэ үр дүн 2026.10.30-ны үйлчилгээнд л хамаарна.
    \item \textbf{Тооны дүрэм.} Цифр ба үгээр бичсэн тоог ялгаатай гэж тооцсон нь Gemini-д хатуу байж болно.
\end{enumerate}

\textbf{Туршилт 2 (хайлт ба хариулт):}
\begin{enumerate}
    \item \textbf{Цөөн асуулт.} 40 асуулт, нэг сайтын хуулбар, нэг өдрийн хэмжилт. Статистикийн ялгааг батлаагүй.
    \item \textbf{Асуултыг хүн хянаагүй.} Асуултыг AI туслах бичиж, сайтын тексттэй тулгаж түгжсэн. Хүн тусад нь хянаагүй.
    \item \textbf{«Утга гарсан» нь бүрэн зөв гэсэн үг биш.} Хариултад шаардсан утга байгаа эсэхийг шалгасан. Хариултыг хүн уншиж үнэлээгүй.
    \item \textbf{Ашиглагдсан багц.} Энэ багцын алдааг харсны дараа систем засвал багц хөгжүүлэлтийн багц болно. Дараагийн шалгалтыг шинэ, хүн хянасан багцаар хийнэ.
    \item \textbf{Бичгийн асуулт.} Дуугаар асуухад яриа танилтын алдаа нэмэгдэнэ. Үүнийг энэ туршилтад хэмжээгүй.
\end{enumerate}
